\documentclass[aspectratio=169]{beamer}

%\documentclass[handout]{beamer}
%% To make 4 per page
%\usepackage{pgfpages}
%\mode<handout>{\setbeamercolor{background canvas}{bg=white}}
%\pgfpagesuselayout{4 on 1}[letterpaper,landscape]%,border shrink=5mm]

\usetheme{default}
% Slide setup, colour independent

\usepackage{amsmath,amssymb,amsthm}
\usepackage{colortbl}
\usepackage{bm}
\usepackage{xcolor}
\usepackage{dsfont}
\usepackage{setspace}
%\usepackage{subfigure}
% To use \ding{234} and the like
\usepackage{pifont}
% To cross reference between slide files
\usepackage{zref-xr,zref-user}
% Use something like
% \zexternaldocument{fileI}
% in the tex files. And cite using \zref instead of \ref

% Fields and the like
\def\IC{\mathbb{C}}
\def\IF{\mathbb{F}}
\def\II{\mathbb{I}}
\def\IJ{\mathbb{J}}
\def\IM{\mathbb{M}}
\def\IN{\mathbb{N}}
\def\IP{\mathbb{P}}
\def\IR{\mathbb{R}}
\def\IZ{\mathbb{Z}}
\def\11{\mathds{1}}


% Bold lowercase
\def\ba{\mathbf{a}}
\def\bb{\mathbf{b}}
\def\bc{\mathbf{c}}
\def\bd{\mathbf{d}}
\def\be{\mathbf{e}}
\def\bf{\mathbf{f}}
\def\bh{\mathbf{h}}
\def\bi{\mathbf{i}}
\def\bj{\mathbf{j}}
\def\bk{\mathbf{k}}
\def\bn{\mathbf{n}}
\def\bp{\mathbf{p}}
\def\br{\mathbf{r}}
\def\bs{\mathbf{s}}
\def\bu{\mathbf{u}}
\def\bv{\mathbf{v}}
\def\bw{\mathbf{w}}
\def\bx{\mathbf{x}}
\def\by{\mathbf{y}}
\def\bz{\mathbf{z}}

% Bold capitals
\def\bB{\mathbf{B}}
\def\bD{\mathbf{D}}
\def\bF{\mathbf{F}}
\def\bG{\mathbf{G}}
\def\bI{\mathbf{I}}
\def\bL{\mathbf{L}}
\def\bN{\mathbf{N}}
\def\bP{\mathbf{P}}
\def\bR{\mathbf{R}}
\def\bS{\mathbf{S}}
\def\bT{\mathbf{T}}
\def\bX{\mathbf{X}}

% Bold numbers
\def\b0{\mathbf{0}}

% Bold greek
\bmdefine{\bmu}{\bm{\mu}}
\def\bphi{\bm{\phi}}
\def\bvarphi{\bm{\varphi}}

% Bold red sentence
\def\boldred#1{{\color{red}\textbf{#1}}}
\def\defword#1{{\color{orange}\textbf{#1}}}

% Caligraphic letters
\def\A{\mathcal{A}}
\def\B{\mathcal{B}}
\def\C{\mathcal{C}}
\def\D{\mathcal{D}}
\def\E{\mathcal{E}}
\def\F{\mathcal{F}}
\def\G{\mathcal{G}}
\def\H{\mathcal{H}}
\def\I{\mathcal{I}}
\def\L{\mathcal{L}}
\def\M{\mathcal{M}}
\def\N{\mathcal{N}}
\def\P{\mathcal{P}}
\def\R{\mathcal{R}}
\def\S{\mathcal{S}}
\def\T{\mathcal{T}}
\def\U{\mathcal{U}}
\def\V{\mathcal{V}}

% tt font for code
\def\code#1{{\tt #1}}

% i.e., e.g.
\def\eg{\emph{e.g.}}
\def\ie{\emph{i.e.}}


% Operators and special symbols
\def\nbOne{{\mathchoice {\rm 1\mskip-4mu l} {\rm 1\mskip-4mu l}
{\rm 1\mskip-4.5mu l} {\rm 1\mskip-5mu l}}}
\def\cov{\ensuremath{\mathsf{cov}}}
\def\Var{\ensuremath{\mathsf{Var}\ }}
\def\Im{\textrm{Im}\;}
\def\Re{\textrm{Re}\;}
\def\det{\ensuremath{\mathsf{det}}}
\def\diag{\ensuremath{\mathsf{diag}}}
\def\nullspace{\ensuremath{\mathsf{null}}}
\def\nullity{\ensuremath{\mathsf{nullity}}}
\def\rank{\ensuremath{\mathsf{rank}}}
\def\range{\ensuremath{\mathsf{range}}}
\def\sgn{\ensuremath{\mathsf{sgn}}}
\def\Span{\ensuremath{\mathsf{span}}}
\def\tr{\ensuremath{\mathsf{tr}}}
\def\imply{$\Rightarrow$}
\def\restrictTo#1#2{\left.#1\right|_{#2}}
\newcommand{\parallelsum}{\mathbin{\!/\mkern-5mu/\!}}
\def\dsum{\mathop{\displaystyle \sum }}%
\def\dind#1#2{_{\substack{#1\\ #2}}}

\DeclareMathOperator{\GL}{GL}
\DeclareMathOperator{\Rel}{Re}
\def\Nt#1{\left|\!\left|\!\left|#1\right|\!\right|\!\right|}
\newcommand{\tripbar}{|\! |\! |}



% The beamer bullet (in base colour)
\def\bbullet{\leavevmode\usebeamertemplate{itemize item}\ }

% Theorems and the like
\newtheorem{proposition}[theorem]{Proposition}
\newtheorem{property}[theorem]{Property}
\newtheorem{importantproperty}[theorem]{Property}
\newtheorem{importanttheorem}[theorem]{Theorem}
%\newtheorem{lemma}[theorem]{Lemma}
%\newtheorem{corollary}[theorem]{Corollary}
\newtheorem{remark}[theorem]{Remark}
\setbeamertemplate{theorems}[numbered]
%\setbeamertemplate{theorems}[ams style]

%
%\usecolortheme{orchid}
%\usecolortheme{orchid}

\def\red{\color[rgb]{1,0,0}}
\def\blue{\color[rgb]{0,0,1}}
\def\green{\color[rgb]{0,1,0}}


% Get rid of navigation stuff
\setbeamertemplate{navigation symbols}{}

% Set footline/header line
\setbeamertemplate{footline}
{%
\quad p. \insertpagenumber \quad--\quad \insertsection\vskip2pt
}
% \setbeamertemplate{headline}
% {%
% \quad\insertsection\hfill p. \insertpagenumber\quad\mbox{}\vskip2pt
% }


\makeatletter
\newlength\beamerleftmargin
\setlength\beamerleftmargin{\Gm@lmargin}
\makeatother

% Colours for special pages
\def\extraContent{yellow!20}


%%%%%%%%%%%%%%%%%
\usepackage{tikz}
\usetikzlibrary{shapes,arrows}
\usetikzlibrary{positioning}
\usetikzlibrary{shapes.symbols,shapes.callouts,patterns}
\usetikzlibrary{calc,fit}
\usetikzlibrary{backgrounds}
\usetikzlibrary{decorations.pathmorphing,fit,petri}
\usetikzlibrary{automata}
\usetikzlibrary{fadings}
\usetikzlibrary{patterns,hobby}

\usetikzlibrary{backgrounds,fit,petri}


\usepackage{pgfplots}
\pgfplotsset{compat=1.6}
\pgfplotsset{ticks=none}

\usetikzlibrary{decorations.markings}
\usetikzlibrary{arrows.meta}
\tikzset{>=stealth}

% For tikz
\usetikzlibrary{shapes,arrows}
\usetikzlibrary{positioning}
\tikzstyle{cloud} = [draw, ellipse,fill=red!20, node distance=0.87cm,
minimum height=2em]
\tikzstyle{line} = [draw, -latex']


%%% For max frame images
\newenvironment{changemargin}[2]{%
\begin{list}{}{%
\setlength{\topsep}{0pt}%
\setlength{\leftmargin}{#1}%
\setlength{\rightmargin}{#2}%
\setlength{\listparindent}{\parindent}%
\setlength{\itemindent}{\parindent}%
\setlength{\parsep}{\parskip}%
}%
\item[]}{\end{list}}


% Make one image take up the entire slide content area in beamer,.:
% centered/centred full-screen image, with title:
% This uses the whole screen except for the 1cm border around it
% all. 128x96mm
\newcommand{\titledFrameImage}[2]{
\begin{frame}{#1}
%\begin{changemargin}{-1cm}{-1cm}
\begin{center}
\includegraphics[width=108mm,height=\textheight,keepaspectratio]{#2}
\end{center}
%\end{changemargin}
\end{frame}
}

% Make one image take up the entire slide content area in beamer.:
% centered/centred full-screen image, no title:
% This uses the whole screen except for the 1cm border around it
% all. 128x96mm
\newcommand{\plainFrameImage}[1]{
\begin{frame}[plain]
%\begin{changemargin}{-1cm}{-1cm}
\begin{center}
\includegraphics[width=108mm,height=76mm,keepaspectratio]{#1}
\end{center}
%\end{changemargin}
\end{frame}
}

% Make one image take up the entire slide area, including borders, in beamer.:
% centered/centred full-screen image, no title:
% This uses the entire whole screen
\newcommand{\maxFrameImage}[1]{
\begin{frame}[plain]
\begin{changemargin}{-1cm}{-1cm}
\begin{center}
\includegraphics[width=\paperwidth,height=\paperheight,keepaspectratio]
{#1}
\end{center}
\end{changemargin}
\end{frame}
}

% This uses the entire whole screen (to include in frame)
\newcommand{\maxFrameImageNoFrame}[1]{
\begin{changemargin}{-1cm}{-1cm}
\begin{center}
\includegraphics[width=\paperwidth,height=0.99\paperheight,keepaspectratio]
{#1}
\end{center}
\end{changemargin}
}

% Make one image take up the entire slide area, including borders, in beamer.:
% centered/centred full-screen image, no title:
% This uses the entire whole screen
\newcommand{\maxFrameImageColor}[2]{
\begin{frame}[plain]
\setbeamercolor{normal text}{bg=#2!20}
\begin{changemargin}{-1cm}{-1cm}
\begin{center}
\includegraphics[width=\paperwidth,height=\paperheight,keepaspectratio]
{#1}
\end{center}
\end{changemargin}
\end{frame}
}


\usepackage{tikz}
\usetikzlibrary{patterns,hobby}
\usepackage{pgfplots}
\pgfplotsset{compat=1.6}
\pgfplotsset{ticks=none}

\usetikzlibrary{backgrounds}
\usetikzlibrary{decorations.markings}
\usetikzlibrary{arrows.meta}
\tikzset{>=stealth}

\tikzset{
  clockwise arrows/.style={
    postaction={
      decorate,
      decoration={
        markings,
        mark=between positions 0.1 and 0.9 step 40pt with {\arrow{>}},
   }}}}


   %%%%%%%%%%%
% To have links to parts in the outline
\makeatletter
\AtBeginPart{%
  \addtocontents{toc}{\protect\beamer@partintoc{\the\c@part}{\beamer@partnameshort}{\the\c@page}}%
}
%% number, shortname, page.
\providecommand\beamer@partintoc[3]{%
  \ifnum\c@tocdepth=-1\relax
    % requesting onlyparts.
    \makebox[6em]{Part #1:} \textcolor{green!30!blue}{\hyperlink{#2}{#2}}
    \par
  \fi
}
\define@key{beamertoc}{onlyparts}[]{%
  \c@tocdepth=-1\relax
}
\makeatother%

\newcommand{\nameofthepart}{}
\newcommand{\nupart}[1]%
    {   \part{#1}%
        \renewcommand{\nameofthepart}{#1}%
        {
          \setbeamercolor{background canvas}{bg=orange!50}
          \begin{frame}{#1}%\partpage 
          \hypertarget{\nameofthepart}{}\tableofcontents%
          \end{frame}
        }
    }



\usecolortheme{orchid}

%% Listings
\usepackage{listings}
\definecolor{mygreen}{rgb}{0,0.6,0}
\definecolor{mygray}{rgb}{0.5,0.5,0.5}
\definecolor{mymauve}{rgb}{0.58,0,0.82}
\definecolor{mygold}{rgb}{1,0.843,0}
\definecolor{myblue}{rgb}{0.537,0.812,0.941}

\definecolor{lgreen}{rgb}{0.6,0.9,.6}
\definecolor{lred}{rgb}{1,0.5,.5}

\lstloadlanguages{R}
\lstset{ %
  language=R,
  backgroundcolor=\color{black!05},   % choose the background color
  basicstyle=\footnotesize\ttfamily,        % size of fonts used for the code
  breaklines=true,                 % automatic line breaking only at whitespace
  captionpos=b,                    % sets the caption-position to bottom
  commentstyle=\color{mygreen},    % comment style
  escapeinside={\%*}{*)},          % if you want to add LaTeX within your code
  keywordstyle=\color{red},       % keyword style
  stringstyle=\color{mygold},     % string literal style
  keepspaces=true,
  columns=fullflexible,
  tabsize=4,
}
% Could also do (in lstset)
% basicstyle==\fontfamily{pcr}\footnotesize
\lstdefinelanguage{Renhanced}%
  {keywords={abbreviate,abline,abs,acos,acosh,action,add1,add,%
      aggregate,alias,Alias,alist,all,anova,any,aov,aperm,append,apply,%
      approx,approxfun,apropos,Arg,args,array,arrows,as,asin,asinh,%
      atan,atan2,atanh,attach,attr,attributes,autoload,autoloader,ave,%
      axis,backsolve,barplot,basename,besselI,besselJ,besselK,besselY,%
      beta,binomial,body,box,boxplot,break,browser,bug,builtins,bxp,by,%
      c,C,call,Call,case,cat,category,cbind,ceiling,character,char,%
      charmatch,check,chol,chol2inv,choose,chull,class,close,cm,codes,%
      coef,coefficients,co,col,colnames,colors,colours,commandArgs,%
      comment,complete,complex,conflicts,Conj,contents,contour,%
      contrasts,contr,control,helmert,contrib,convolve,cooks,coords,%
      distance,coplot,cor,cos,cosh,count,fields,cov,covratio,wt,CRAN,%
      create,crossprod,cummax,cummin,cumprod,cumsum,curve,cut,cycle,D,%
      data,dataentry,date,dbeta,dbinom,dcauchy,dchisq,de,debug,%
      debugger,Defunct,default,delay,delete,deltat,demo,de,density,%
      deparse,dependencies,Deprecated,deriv,description,detach,%
      dev2bitmap,dev,cur,deviance,off,prev,,dexp,df,dfbetas,dffits,%
      dgamma,dgeom,dget,dhyper,diag,diff,digamma,dim,dimnames,dir,%
      dirname,dlnorm,dlogis,dnbinom,dnchisq,dnorm,do,dotplot,double,%
      download,dpois,dput,drop,drop1,dsignrank,dt,dummy,dump,dunif,%
      duplicated,dweibull,dwilcox,dyn,edit,eff,effects,eigen,else,%
      emacs,end,environment,env,erase,eval,equal,evalq,example,exists,%
      exit,exp,expand,expression,External,extract,extractAIC,factor,%
      fail,family,fft,file,filled,find,fitted,fivenum,fix,floor,for,%
      For,formals,format,formatC,formula,Fortran,forwardsolve,frame,%
      frequency,ftable,ftable2table,function,gamma,Gamma,gammaCody,%
      gaussian,gc,gcinfo,gctorture,get,getenv,geterrmessage,getOption,%
      getwd,gl,glm,globalenv,gnome,GNOME,graphics,gray,grep,grey,grid,%
      gsub,hasTsp,hat,heat,help,hist,home,hsv,httpclient,I,identify,if,%
      ifelse,Im,image,\%in\%,index,influence,measures,inherits,install,%
      installed,integer,interaction,interactive,Internal,intersect,%
      inverse,invisible,IQR,is,jitter,kappa,kronecker,labels,lapply,%
      layout,lbeta,lchoose,lcm,legend,length,levels,lgamma,library,%
      licence,license,lines,list,lm,load,local,locator,log,log10,log1p,%
      log2,logical,loglin,lower,lowess,ls,lsfit,lsf,ls,machine,Machine,%
      mad,mahalanobis,make,link,margin,match,Math,matlines,mat,matplot,%
      matpoints,matrix,max,mean,median,memory,menu,merge,methods,min,%
      missing,Mod,mode,model,response,mosaicplot,mtext,mvfft,na,nan,%
      names,omit,nargs,nchar,ncol,NCOL,new,next,NextMethod,nextn,%
      nlevels,nlm,noquote,NotYetImplemented,NotYetUsed,nrow,NROW,null,%
      numeric,\%o\%,objects,offset,old,on,Ops,optim,optimise,optimize,%
      options,or,order,ordered,outer,package,packages,page,pairlist,%
      pairs,palette,panel,par,parent,parse,paste,path,pbeta,pbinom,%
      pcauchy,pchisq,pentagamma,persp,pexp,pf,pgamma,pgeom,phyper,pico,%
      pictex,piechart,Platform,plnorm,plogis,plot,pmatch,pmax,pmin,%
      pnbinom,pnchisq,pnorm,points,poisson,poly,polygon,polyroot,pos,%
      postscript,power,ppoints,ppois,predict,preplot,pretty,Primitive,%
      print,prmatrix,proc,prod,profile,proj,prompt,prop,provide,%
      psignrank,ps,pt,ptukey,punif,pweibull,pwilcox,q,qbeta,qbinom,%
      qcauchy,qchisq,qexp,qf,qgamma,qgeom,qhyper,qlnorm,qlogis,qnbinom,%
      qnchisq,qnorm,qpois,qqline,qqnorm,qqplot,qr,Q,qty,qy,qsignrank,%
      qt,qtukey,quantile,quasi,quit,qunif,quote,qweibull,qwilcox,%
      rainbow,range,rank,rbeta,rbind,rbinom,rcauchy,rchisq,Re,read,csv,%
      csv2,fwf,readline,socket,real,Recall,rect,reformulate,regexpr,%
      relevel,remove,rep,repeat,replace,replications,report,require,%
      resid,residuals,restart,return,rev,rexp,rf,rgamma,rgb,rgeom,R,%
      rhyper,rle,rlnorm,rlogis,rm,rnbinom,RNGkind,rnorm,round,row,%
      rownames,rowsum,rpois,rsignrank,rstandard,rstudent,rt,rug,runif,%
      rweibull,rwilcox,sample,sapply,save,scale,scan,scan,screen,sd,se,%
      search,searchpaths,segments,seq,sequence,setdiff,setequal,set,%
      setwd,show,sign,signif,sin,single,sinh,sink,solve,sort,source,%
      spline,splinefun,split,sqrt,stars,start,stat,stem,step,stop,%
      storage,strstrheight,stripplot,strsplit,structure,strwidth,sub,%
      subset,substitute,substr,substring,sum,summary,sunflowerplot,svd,%
      sweep,switch,symbol,symbols,symnum,sys,status,system,t,table,%
      tabulate,tan,tanh,tapply,tempfile,terms,terrain,tetragamma,text,%
      time,title,topo,trace,traceback,transform,tri,trigamma,trunc,try,%
      ts,tsp,typeof,unclass,undebug,undoc,union,unique,uniroot,unix,%
      unlink,unlist,unname,untrace,update,upper,url,UseMethod,var,%
      variable,vector,Version,vi,warning,warnings,weighted,weights,%
      which,while,window,write,\%x\%,x11,X11,xedit,xemacs,xinch,xor,%
      xpdrows,xy,xyinch,yinch,zapsmall,zip},%
   otherkeywords={!,!=,~,$,*,\%,\&,\%/\%,\%*\%,\%\%,<-,<<-,_,/},%
   alsoother={._$},%
   sensitive,%
   morecomment=[l]\#,%
   morestring=[d]",%
   morestring=[d]'% 2001 Robert Denham
  }%

%%%%%%% 
%% Definitions in yellow boxes
\usepackage{etoolbox}
\setbeamercolor{block title}{use=structure,fg=structure.fg,bg=structure.fg!40!bg}
\setbeamercolor{block body}{parent=normal text,use=block title,bg=block title.bg!20!bg}

\BeforeBeginEnvironment{definition}{%
	\setbeamercolor{block title}{fg=black,bg=yellow!20!white}
	\setbeamercolor{block body}{fg=black, bg=yellow!05!white}
}
\AfterEndEnvironment{definition}{
	\setbeamercolor{block title}{use=structure,fg=structure.fg,bg=structure.fg!20!bg}
	\setbeamercolor{block body}{parent=normal text,use=block title,bg=block title.bg!50!bg, fg=black}
}
\BeforeBeginEnvironment{importanttheorem}{%
	\setbeamercolor{block title}{fg=black,bg=red!20!white}
	\setbeamercolor{block body}{fg=black, bg=red!05!white}
}
\AfterEndEnvironment{importanttheorem}{
	\setbeamercolor{block title}{use=structure,fg=structure.fg,bg=structure.fg!20!bg}
	\setbeamercolor{block body}{parent=normal text,use=block title,bg=block title.bg!50!bg, fg=black}
}
\BeforeBeginEnvironment{importantproperty}{%
	\setbeamercolor{block title}{fg=black,bg=red!50!white}
	\setbeamercolor{block body}{fg=black, bg=red!30!white}
}
\AfterEndEnvironment{importantproperty}{
	\setbeamercolor{block title}{use=structure,fg=structure.fg,bg=structure.fg!20!bg}
	\setbeamercolor{block body}{parent=normal text,use=block title,bg=block title.bg!50!bg, fg=black}
}


% Beginning of a section
\AtBeginSection[]{
	{
		\setbeamercolor{background canvas}{bg=yellow!10}
		\begin{frame}[noframenumbering,plain]
			\framesubtitle{\nameofthepart Chapter \insertromanpartnumber \ -- \iteminsert{\insertpart}}
			\tableofcontents[
				currentsection,
				sectionstyle=show/shaded,
				subsectionstyle=show/show/hide]
		\end{frame}
	\addtocounter{page}{-1}
	%\addtocounter{framenumber}{-1} 
	}
}

% Beginning of a section
\AtBeginSubsection[]{
	{
		\setbeamercolor{background canvas}{bg=green!10}
		\begin{frame}[noframenumbering,plain]
				\framesubtitle{\nameofthepart Chapter \insertromanpartnumber \ -- \iteminsert{\insertpart}}
				\tableofcontents[
					currentsection,
					sectionstyle=show/show,
					currentsubsection,
					subsectionstyle=show/shaded/hide]
			\end{frame}
		\addtocounter{page}{-1}
	}
}


% To have theorems and everything derived prefixed by the slide set
% number
\renewcommand{\thetheorem}{6.\arabic{theorem}}



\title{Nonnegative matrices}
\author{Julien Arino}
\date{Fall 2023}

%%%%%%%%%%%%%%%%%%%%%
%%%%%%%%%%%%%%%%%%%%%
%%%%%%%%%%%%%%%%%%%%%
%%%%%%%%%%%%%%%%%%%%%
%%%%%%%%%%%%%%%%%%%%%
%%%%%%%%%%%%%%%%%%%%%
%%%%%%%%%%%%%%%%%%%%%
%%%%%%%%%%%%%%%%%%%%%
\begin{document}

% The title page
\begin{frame}[noframenumbering,plain]
  \begin{tikzpicture}[remember picture,overlay]
    \node[above right,inner sep=0pt] at (current page.south west)
    {
        \includegraphics[width=\paperwidth]{title-page-picture.png}
    };
\end{tikzpicture}
	\titlepage
\end{frame}
\addtocounter{page}{-1}
  
  
\begin{frame}{Outline}
	  \tableofcontents[hideallsubsections]
\end{frame}
\addtocounter{page}{-1}


%%%%%%%%%%%%%%%%%%%%%%%%%%%
%%%%%%%%%%%%%%%%%%%%%%%%%%%
%%%%%%%%%%%%%%%%%%%%%%%%%%%
%%%%%%%%%%%%%%%%%%%%%%%%%%%
\section{Definitions and some preliminary results}

\begin{frame}
\begin{definition}[Nonnegative/positive matrix]
  \label{def:nonnegative_matrix}
A matrix $A\in \M_{mn}(\IR)$ is a \defword{nonnegative matrix} if $a_{ij}\geq 0$ for all $i=1,\ldots,m$ and $j=1,\ldots,n$. We write $A\geq 0$. $A$ is a \defword{positive matrix} if $a_{ij}>0$ for all $i=1,\ldots,m$ and $j=1,\ldots,n$. We write $A>0$
\end{definition}
\vfill
\begin{remark}
In other references, you will see
\begin{itemize}
    \item $A\geq 0 \iff a_{ij}\geq 0$
    \item $A>0 \iff A\geq 0$ and there exists $(i,j)$, $a_{ij}>0$ 
    \hfill [positive]
    \item $A\gg 0 \iff a_{ij}>0$ for all $i,j$ \hfill [strongly positive]
\end{itemize}
I tend to favour the latter notation over the one used in these notes, but since the former is more common in matrix theory, I use the notation of Definition~\ref{def:nonnegative_matrix} here
\end{remark}
\end{frame}


\begin{frame}{Notation}
Let $A,B\in\M_{mn}(\IR)$. Nonnegativity and positivity are used to define partial orders on $\M_{mn}(\IR)$
\begin{itemize}
    \item $A\geq B \iff A-B\geq 0$
    \item $A>B \iff A-B >0$
\end{itemize}
The same is used for vectors $\bx,\by\in\IR^n$: $\bx\geq\by$ and $\bx>\by$ if, respectively, $\bx-\by\geq\b0$ and $\bx-\by>\b0$. Note that the order is only partial: if $A\geq 0$ and $B\geq 0$, for instance, it is not necessarily possible to decide on the ordering of $A$ and $B$ with respect to one another
\end{frame}


\begin{frame}
\begin{theorem}\label{th:nonnegative_matrices_semigroup}
Let $A$ and $B$ be nonnegative matrices of appropriate sizes. Then $A+B$ and $AB$ are nonnegative. 
If $A>0$ and $B\geq 0$, $B\neq 0$, then $AB\geq 0$ and $AB\neq 0$
\end{theorem}
\vfill
\begin{corollary}\label{coro:nonnegativity_preserves_orientation}
Let $\bx,\by\in\IR^n$ be such that $\bx\geq\by$ and $A\in \M_{mn}$ be nonnegative. Then $A\bx\geq A\by$. 
Assume additionally that $\bx\geq\by$, $\bx\neq\by$ and $A>0$. Then $A\bx>A\by$
\end{corollary}
\end{frame}


%%%%%%%%%%%%%%%%%%%%%%%%
%%%%%%%%%%%%%%%%%%%%%%%%
\subsection{Zero-nonzero structure of a matrix}
\begin{frame}
\begin{definition}
Let $P,Q\in \M_{nm}(\IF)$. $P$ and $Q$ have the same \defword{zero-nonzero structure} if for all $i,j$, $p_{ij}\neq 0 \iff q_{ij}\neq 0$
\end{definition}
\vfill
Zero-nonzero structure defines an equivalence relation. Therefore, as with all equivalence relations, one only needs one representative from the equivalence class. One typical representative is defined using Boolean matrices
\end{frame}


\begin{frame}
\begin{definition}
A \defword{Boolean matrix}\index{Boolean matrix} is a matrix whose entries are Boolean $\{0,1\}$ and use Boolean arithmetics:
\begin{itemize}
    \item $0+0=0$
    \item $1+0=0+1=1$
    \item $1+1=1$
    \item $0\cdot 1= 1$ and $ 1= 0= 0\cdot 0$
    \item $1\cdot 1=1$
\end{itemize}
\end{definition}
\begin{definition} 
Let $A\in \M_{nm}(\IF)$.
Then $A_B$ denotes the \defword{Boolean representation} of $A$, defined as follows. If $A=[a_{ij}]$, then $A_B=[\alpha_{ij}]$ with
\[
\alpha_{ij} = 
\begin{cases}
1& \textrm{if }a_{ij}\neq 0\\
0& \textrm{if }a_{ij}=0
\end{cases}
\]
\end{definition}

\begin{theorem}
Let $P,], Q\in \M_{n}$ be nonnegative matrices. The zero-nonzero structure of the sum of product of $P$ and $Q$ is uniquely determined by the zero-nonzero structure of $P$ and $Q$. Also,
\begin{align*}
    (P+B)_B &= P_B+Q_B\\
    (PQ)_B& = P_BQ_B
\end{align*}
 (with Boolean addition and multiplication on the right).
\end{theorem}

\begin{theorem}\label{th:aij_positive_path_length_k}
Let $A\in \M_n$ be a nonnegative matrix, $k\in \IN\setminus\{0\}$. Then the $(i,j)$ entry in $A^k$ is nonzero if and only if there is a directed walk of length \emph{exactly} $k$ from $i$ to $j$ in $G(A)$, the digraph associated to $A$.
\end{theorem}
\begin{remark}
The walk can go thought the same vertex more than once.
\end{remark}

\begin{theorem}\label{th:sum_of_irreducible_matrices_positive}
Let $A\in \M_n$ be a nonnegative and irreducible, $\alpha_0, \alpha_1, \ldots, \alpha_{n-1}\in \IR_{+}\setminus\{0\}$. Then 
\[
\alpha_0\II+\alpha_1A+\cdots+ \alpha_{n-1}A^{n-1}>0.
\]
In particular, $(\II+A)^{n-1}>0$.
\end{theorem}


\begin{definition}
Let $A\in \M_n(\IF)$. The matrix $m(A)= [| a_{ij}|]$ is the \defword{modulus}\index{matrix!modulus} of $A$.
\end{definition}
\begin{lemma}
$A,\, B\in \M_n(\IF)$, $A$ potentially complex. Suppose $m(A)<B$. Then $\rho(A)\leq \rho(B)$. In particular, $\rho(A)\leq \rho(m(A))$.
\end{lemma}

\begin{lemma}\label{AnonnegativeB}
Let $A\in \M_n$ be a nonnegative matrix, $0 \neq z\in \IR^n  $. If $\exists b\in \IR$ such that $Az>bz$, then $\rho(A)>b$.
\end{lemma}
Note that the latter result also holds true if $>$ is replaced by $\geq$. See \cite[Theorem 8.3.2]{HornJohnson2013} and a proof there in the latter case.
Note, however, that the proof of \cite[Theorem 8.3.2]{HornJohnson2013} uses Perron's Lemma. So it is better for our purpose here to use the proof in \cite{Fiedler2008}, where this result is Lemma 4.1.6.
\end{frame}

%%%%%%%%%%%%%%%%%%%%%%%%
%%%%%%%%%%%%%%%%%%%%%%%%
%%%%%%%%%%%%%%%%%%%%%%%%
%%%%%%%%%%%%%%%%%%%%%%%%
\section{The Perron-Frobenius theorem}
\label{sec:PerronFrobenis}



%%%%%%%%%%%%%%%%%%%%%%%%
%%%%%%%%%%%%%%%%%%%%%%%%
\subsection{The Perron-Frobenius Theorem for irreducible matrices}
\label{sec:PF_irreducible_matrices}

\begin{frame}
\begin{theorem}[Perron-Frobenius]\index{Perron-Frobenius Theorem}\label{th:PF}
Let $A\geq 0\in \M_n$ be irreducible. Then the spectral radius $\rho(A)= \max \{ |\lambda|, \lambda\in \sigma(A)\}$ is an eigenvalue of $A$. It is simple (has algebraic multiplicity 1), positive and is associated with a positive eigenvector. Furthermore, there is no nonnegative eigenvector associated to any other eigenvalue of $A$
\end{theorem}
\vfill
\begin{remark}
We often say that $\rho(A)$ is the \defword{Perron root} of $A$; the corresponding eigenvector is the \defword{Perron vector} of $A$
\end{remark}
\end{frame}


%%%%%%%%%%%%%%%%%%%%%%%%
%%%%%%%%%%%%%%%%%%%%%%%%
\subsection{Proof of a result of Perron for positive matrices}
\label{sec:proof_Perron}

\begin{frame}
\begin{lemma}[Perron]\index{Perron's Lemma}\label{lemma:Perron}
Let $\M_n\ni A>0$. Then $\rho(A)$ is a positive eigenvalue of $A$ and there is only one linearly independent eigenvector associated to $\rho(A)$, which can be taken to be positive
\end{lemma}
\vfill
\begin{lemma}\label{lemma:prep_Perron}
Let $\alpha_1,\ldots,\alpha_n\in\IR^{*}_{+}$ and $v_1,\ldots,v_n\in\IC$. Then
\begin{equation}\label{ineq:proof_Perron_1}
\left|\sum\limits_{i=1}^{n} \alpha_i v_i\right| 
\leq 
\sum\limits_{i=1}^n \alpha |v_i|
\end{equation}
with equality if and only if there exists $\eta\in\IC$, $|\eta|=1$, such that $\eta v_i\geq 0$ for all $i=1, \ldots, n$
\end{lemma}
\end{frame}

%%%%%%%%%%%%%%%%%%%%%%%%
%%%%%%%%%%%%%%%%%%%%%%%%
\subsection{Proof of the Perron-Frobenius theorem for irreducible matrices}
\label{sec:proof_PF}


\begin{frame}
\begin{theorem}\label{CoolTh1}
Let $A\in \M_{n}$ and $f(x)$ a polynomial. Then 
\[
\sigma(f(A))= \{ f(\lambda_1), \ldots, f(\lambda_n), \lambda_i\in \sigma(A)\}
\]
If we have $g(\lambda_i)\neq 0$ for $\lambda_i\in \sigma(A)$, for some polynomial $g$, then the matrix $g(A)$ is non-singular and 
\[
\sigma\left(f(A)g(A)^{-1}\right)
= 
\left\lbrace \dfrac{f(\lambda_1)}{g(\lambda_1)},\ldots, \dfrac{f(\lambda_n)}{g(\lambda_n)}, \lambda_i\in \sigma(A) \right\rbrace
\]
If $x\neq 0$ eigenvector of $A$ associated to $\lambda\in \sigma(A)$, then $x$ is also an eigenvector of $f(A)$ and $f(A)g(A)^{-1}$ associated to eigenvalue $f(\lambda)$ and $f(\lambda)/g(\lambda)$, respectively
\end{theorem}
\end{frame}


\begin{frame}
\begin{lemma}[Schur's lemma]\label{lemma:Schur_LI_evectors}
Let $A\in \M_n$ and $\lambda\in \sigma(A)$. 
Then $\lambda$ is simple if and only if both the following conditions are statisfied:
\begin{enumerate}
     \item There exists only one linear independent eigenvector of $A$ associated to $\lambda$, say $\bu$, and thus only one linear independent eigenvector of $A^T$ associated to $\lambda$, say $\bv$
     \item Vectors $\bu$ and $\bv$ in (1) satisfy $\bv^T\bu\neq 0$
\end{enumerate}
\end{lemma}
\end{frame}


%%%%%%%%%%%%%%%%%%
%%%%%%%%%%%%%%%%%%
\subsection{Primitive matrices}
\label{sec:primitivity}


\begin{frame}
\begin{definition}
Let $\M_n(\IR)\ni A\geq 0$.
We say that $A$ is \defword{primitive} (with \defword{primitivity index} $k\in \IN_{+}^{*}$) if there exists $k\in \IN_{+}^{*}$ such that 
\[
A^k>0
\]
with $k$ the smallest integer for which this is true.
We say that a matrix is \defword{imprimitive} if it is not primitive
\end{definition}
\vfill
\begin{remark}
Primitivity implies irreducibility. The converse is not true
\end{remark}
\end{frame}


\begin{frame}
\begin{theorem}
A sufficient condition for primitivity is irreducibility with at least one positive diagonal entry
\end{theorem}
\vfill
Here $d$ is the index of imprimitivity (i.e., the number of eigenvalues that have the same modulus as $\lambda_p=\rho(A)$). If $d=1$, then $A$ is primitive. We have that $d=\gcd$ of all the lengths of closed walks in $G(A)$
\vfill
\begin{theorem}
Let $A\in \M_n$ be a non-negative matrix. If $A$ is primitive, then $A^k>0$ for some $0<k\leq (n-1)n^n$
\end{theorem}
\end{frame}


\begin{frame}
\begin{theorem}
Let $A\geq 0$ primtive. Suppose the shortest simple directed cycle in $G(A)$ has length $s$, then primitivity index is $\leq n+s(n-1)$
\end{theorem}
\vfill
\begin{theorem}
Let $A\in \M_n$ be a nonnegative matrix. $A$ is primitive if and only if $A^{n^2-2n+2}>0$
\end{theorem}
\vfill
\begin{theorem}
Let $A\in \M_n$ be a nonnegative irreducible matrix . Suppose that $A$ has $d$ positive entries on the diagonal. Then the primitivity index is $\leq 2n-d-1$
\end{theorem}
\end{frame}


%%%%%%%%%%%%%%%%%%%%%%%%
%%%%%%%%%%%%%%%%%%%%%%%%
\subsection{The Perron-Frobenius Theorem for nonnegative matrices}
\label{sec:PF_nonnegative_matrices}


\begin{frame}
\begin{theorem}
Let $A\geq 0$ in $\M_n$. Then there exists $0\neq v\geq 0$ such that $A\bv= \rho(A)\bv$
\end{theorem}
\end{frame}

%%%%%%%%%%%%%%%%%%%%%%%%
%%%%%%%%%%%%%%%%%%%%%%%%
\subsection{The Perron-Frobenius Theorem (revamped)}
\label{sec:PF_summary}

\begin{frame}
Let us restate the Perron-Frobenius theorem, taking into account the different cases.
The classification in the following result is inspired by the presentation in \cite{Caswell2001}.
\begin{theorem}\label{th:PF_allCases}
Let $\M_n\ni A\geq 0$. Denote $\lambda_P$ the Perron root of $A$, i.e., $\lambda_P=\rho(A)$, $\bv_P$ and $\bw_P$ the corresponding right and left Perron vectors of $A$, respectively. 
Denote $d$ the index of imprimitivity of $A$ (with $d=1$ when $A$ is primitive) and $\lambda_j\in\sigma(A)$ the spectrum of $A$, with $j=2,\ldots,n$ unless otherwise specified (assuming $\lambda_1=\lambda_P$).
Then conclusions of the Perron-Frobenius Theorem can be summarised as follows.
\end{theorem}
\end{frame}

\begin{frame}
	\centering
	\begin{tikzpicture}[auto, %node distance = 2cm, auto,
	cloud/.style={minimum width={width("N-1")+2pt},
		draw, rectangle,fill=gray!10}]
	\node [cloud] (nonneg) {\sc Nonnegative};
	\node [cloud] (reduc) [below left = of nonneg,text width=3cm] {{\sc Reducible}\\ \noindent\begin{itemize}
		% \renewcommand\labelitemi{\ding{234}}
		\item $\lambda_P\geq 0$
		\item $\bw_P\geq 0$
		\item $\bv_P\geq 0$
		\item $\lambda_P \geq | \lambda_j |$
		\end{itemize}
	};
	\node [cloud] (irred) [below right = of nonneg] {\sc Irreducible};	
	\node [cloud] (notprim) [below left = of irred,text width=3cm] {{\sc Imprimitive}\\ \begin{itemize}
		% \renewcommand\labelitemi{\ding{234}}
		\item $\lambda_P> 0$
		\item $\bw_P> 0$
		\item $\bv_P> 0$
		\item $\lambda_P=|\lambda_j|$, $j=2,\ldots,d$
		\item $\lambda_P > |\lambda_j| $, $j>d$
		\end{itemize}
	};	
	\node [cloud] (prim) [below right = of irred,text width=3cm] {{\sc Primitive}\\
		\begin{itemize}
		% \renewcommand\labelitemi{\ding{234}}
		\item $\lambda_P> 0$
		\item $\bw_P> 0$
		\item $\bv_P> 0$
		\item $\lambda_P>|\lambda_j|$, $j\neq P$ 
		\end{itemize}
	};	
	%% Links
	\path [line, thick] (nonneg) to (reduc);
	\path [line, thick] (nonneg) to (irred);
	\path [line, thick] (irred) to (notprim);
	\path [line, thick] (irred) to (prim);
	\end{tikzpicture}
\end{frame}



%%%%%%%%%%%%%%%%%%%%%%%%
%%%%%%%%%%%%%%%%%%%%%%%%
\subsection{Application of the Perron-Frobenius Theorem}
\label{sec:applications_PF}


\begin{frame}
\begin{theorem}
Let $A\in \M_n$ be a nonnegative irreducible matrix and $\in \IN_{+}$. Then the following ar eequivalent:
\begin{enumerate}
    \item there exists exactly $h$ distinct eigenvalues such that $|\lambda| = \rho(A)$.
    \item there exists $P$ a permutation matrix such that 
    \[PAP^{T}= \begin{pmatrix}
    0& A_{12}&0 & \dots & 0\\
    \vdots& & A_{23}& \ddots&\vdots \\
    \vdots&& &\ddots&0 \\
    \vdots&& && A_{h-1,h}\\
    A_{h1} &0&\ldots& \ldots & 0
    \end{pmatrix}\]
    where the diagonal blocks are square, and there does not exists other permutation matrix giving less than $h$ horizontal blocks.
    \item the greatest common divisor of the lengths of all cycles in $G(A)$ is $h$.
    \item $h$ is the maximal positive integer $k$ such that 
    \[\sigma(e^{2\pi i /k}A)= \sigma(A)\]
\end{enumerate}
\end{theorem}
\end{frame}


\begin{frame}
\begin{corollary}
Let $A\in \M_n$, $A\geq 0$ irreducible with exactly $h$ distinct eigenvalues of modulus $\rho(A)$. Then, we can consider this eigenvalues as points in the complex plan, the eigenvalues are the vertices of a regular polygon of $h$ sides with centre at the origin and are of the vertices being $\rho(A)$

\begin{tikzpicture}
\node  (I0) at (1.4,-0.3) {$\rho(A)$};
\node  (I0) at (1.4,-1.3) {$h=1$};
\draw[thick,->] (-2,0) -- (2,0) node[anchor=north west] {};
\draw[thick,->] (0,-2) -- (0,2) node[anchor=south east] {};
\draw (0,0) circle (1cm);

\node  (I0) at (6.4,-0.3) {$\rho(A)$};
\node  (I1) at (6,0) {$\bullet$};
\node  (I1) at (4,0) {$\bullet$};
\node  (I0) at (6.4,-1.3) {$h=2$};
\draw[thick,->] (3,0) -- (7,0) node[anchor=north west] {};
\draw[thick,->] (5,-2) -- (5,2) node[anchor=south east] {};
\draw (5,0) circle (1cm);

\node  (I0) at (11.4,-0.3) {$\rho(A)$};
\node  (I1) at (9.5,-0.87) {$\bullet$};
\node  (I1) at (9.5,0.87) {$\bullet$};
\node  (I1) at (11,0) {$\bullet$};
\node  (I0) at (11.4,-1.3) {$h=3$};
\draw[thick,->] (8,0) -- (12,0) node[anchor=north west] {};
\draw[thick,->] (10,-2) -- (10,2) node[anchor=south east] {};
\draw (10,0) circle (1cm);
\end{tikzpicture}
\end{corollary}
\end{frame}



\begin{frame}
\begin{remark}
For Fiedler, a primitive matrix is defined as an irreducible nonnegative matrix such that $h=1$
\end{remark}
\end{frame}


\begin{frame}
\begin{theorem}
Let $A\geq 0$ in $\M_n$, $n \geq 2$. TFAE
\begin{enumerate}
    \item $A^n=0$
    \item there exists $\IN\ni k>0$ such that $A^k=0$
    \item $G(A)$ acyclic
    \item $\exists P$, permutation matrix, .t. $PAP^T$ is upper-triangular with zeros on main diagonal
    \item $\rho(A)=0$
\end{enumerate}
\end{theorem}
\vfill
\begin{theorem}
Let $A\geq 0$ be a nonnegative matrix in $\M_n$. Assume that $A$ has a positive eigenvector. Then that eigenvector is the Perron vector and is associated to $\rho(A)$
\end{theorem}
\end{frame}


%%%%%%%%%%%%%%%%%%%%%%%%%%%
%%%%%%%%%%%%%%%%%%%%%%%%%%%
%%%%%%%%%%%%%%%%%%%%%%%%%%%
%%%%%%%%%%%%%%%%%%%%%%%%%%%
\section{Stochastic matrices}

%%%%%%%%%%%%%%%%%%%%%%%%%%%
%%%%%%%%%%%%%%%%%%%%%%%%%%%
\subsection{Row- and column-stochastic matrices}\label{sec:rowStochastic_matrices}


\begin{frame}
\begin{definition}[Stochastic matrix]
The matrix $A\in \M_n$ is \defword{stochastic} if
\begin{enumerate}
    \item[-] $A\geq 0$ \hfill[The matrix is nonnegative]
    \item[-] $A\nbOne = \nbOne$,  $\nbOne=(1,\ldots,1)^T$ \hfill[All rows sum to 1]
\end{enumerate}
\end{definition}
\vfill
Equivalently, the matrix is stochastic if its column sums all equal 1
\vfill
\begin{definition}
  The matrix is \defword{row-stochastic} or \defword{column-stochastic}, respectively, if the rows or columns sum to 1. The terms \defword{right stochastic} and \defword{left stochastic} are also used. If both rows and columns sum to 1, then the matrix is \defword{doubly stochastic}
\end{definition}
\end{frame}


\begin{frame}
\begin{theorem}
Let $A\in\M_n$ be stochastic. Then $\rho(A)=1$
\end{theorem}
\vfill
\begin{theorem}
Let $P\in\M_n$, $P\geq 0$. Assume that $P$ has a positive eigenvector $u$ and that $\rho(P)>0$. Then there exists $D$, diagonal matrix with $\diag(D)>0$, and $k>0$, $k\in \IR$ such that 
\[
A=kDPD^{-1}
\]
is stochastic, with $k=\rho(P)^{-1}$
\end{theorem}
\end{frame}



\begin{frame}
\begin{theorem}
Let $A,\, B\in \M_n$ be stochastic. Then $AB$ is stochastic
\end{theorem}
\vfill
\begin{theorem}
Let $A$ be a primitive stochastic. Then $A^{k}\to \nbOne \bv^T$, $k \to \infty$, where $\nbOne \bv^T$ has rank 1 and $\bv$ is the (left) eigenvector of $A^T$ associated to $\rho(A)=1$ and normalised so that $\bv^T \nbOne=1$
\end{theorem}
\vfill
\begin{remark}
This is a result that is used to compute the limit of a regular Markov chain
\end{remark}
\end{frame}


%%%%%%%%%%%%%%%%%%%%%%%%%%%
%%%%%%%%%%%%%%%%%%%%%%%%%%%
\subsection{Doubly stochastic matrices}\label{sec:doublyStochastic_matrices}

\begin{frame}
\begin{definition}
The matrix $A\in \M_n$, $A\geq 0$ is \defword{doubly stochastic}\index{matrix!doubly stochastic} if $A\nbOne =\nbOne$ and $\nbOne^T A= \nbOne^T$
\end{definition}
\vfill
\begin{remark}
Here $\rho(A)=1$ is associated to $\nbOne$ for $A$ and for $A^T$
\end{remark}
\end{frame}


\begin{frame}
Consider $E$ the Euclidean space. A set $K$ of points in $E$ is \defword{convex} if $A_1, \, A_2$ points in $K$, $\lambda_1, \lambda_2\in \IR_+$ such that $\lambda_1+\lambda_2=1$, then 
\[\lambda_1 A_1+ \lambda_2 A_2\in K.\]
A \defword{convex polyhedron} $K$ is the set of all points of the form 
\[
  \sum\limits_{i=1}^N \lambda_i A_i
  \]
where $A_i$ are points in $E$ and $\lambda_1 \in \IR_+$
\vfill
Let $A=[a_{ij}]\in\M_n(\IR)$. Consider this matrix as a point in $E$ with coordinates $[a_{11}, a_{12}, \dots, a_{nn}]$ ($\dim E= n^2$)
\end{frame}

\begin{frame}
\begin{theorem}
 Let $A\in \M_{n}$, $A=[a_{ij}]$, if $A$ is doubly stochastic, then this forms an $(n-1)^2$ dimensional subspace of $\tilde{E}= \IR^{n^2}$
\end{theorem}
\vfill
\begin{theorem}[Birkhoff]
In the space $\tilde{E}=R^{n^2}$, the set of doubly stochastic matrices of order $n$ is a convex polyhedron in $E$ (the subspace of stochastic matrices). The vertices of the polyhedron are the points corresponding to all the permutation matrices
\end{theorem}
\end{frame}


%%%%%%%%%%%%%%%%%%
%%%%%%%%%%%%%%%%%%
%%%%%%%%%%%%%%%%%%
%%%%%%%%%%%%%%%%%%
\begin{frame}[allowframebreaks]
    \frametitle{References}
    \bibliographystyle{amsalpha}
    \bibliography{MATH-4370-7370-lecture-notes.bib}
\end{frame}
    
\end{document}